\chapter{Universe, Symbology and Corporate Actions}\label{ch:rs:universe-symbology-and-corporate-actions}

Before the market opened on 9 June 2022, the Class A shares of the company formerly called Facebook began trading on
Nasdaq under a new ticker, META; the ticker FB, used since the listing in 2012, was retired. A research table keyed on
tickers either loses ten years of history at that date or, if the old ticker is later given to someone else, splices two
companies into one series. Tickers are not the only thing that moves: companies list and delist, split their shares, merge
and spin off, and the set of stocks a strategy could have traded changes every month. This chapter builds the objects that
keep all of this straight: a universe defined through time, a \omterm{def:rs:universe-symbology-and-corporate-actions:master}{security master} keyed by \omterm{def:rs:universe-symbology-and-corporate-actions:permanent}{permanent identifiers}, adjustments
that do not leak the future, and \omterm{def:rs:universe-symbology-and-corporate-actions:delisting}{delisting returns}. Its data are a simulated market of about 1\,500 listings over ten years,
in which tickers change and are reused, shares split and companies fail or are bought.

\section{A tradable universe through time}

\begin{definition}[Tradable universe, universe membership, liquidity filter]\label{def:rs:universe-symbology-and-corporate-actions:universe}
A \emph{tradable universe}\index{tradable universe} is the set of securities a strategy may hold at a given date, defined by
rules evaluated on data known at that date. \emph{Universe membership}\index{universe membership} is the resulting table of
(date, security) pairs. A \emph{liquidity filter}\index{liquidity filter} is a membership rule on tradability: a minimum
average daily traded value, a minimum price, a minimum free-float capitalisation, a rank by any of these.
\end{definition}

A universe is a research decision with consequences. It fixes which names the predictors are measured on (chapter~6), which
the optimiser may buy (chapter~25) and how much capital the strategy can take (chapter~28). Two properties make it honest:
its rules use only data known at each date, and it does not change more than the underlying liquidity does. The second is
not automatic. A rule that takes the 500 names with the largest traded value each month replaces every name whose rank
wanders across 500, although nothing about the name has changed; each replacement is a pair of trades the strategy then
pays for.

\begin{method}[A buffered membership rule]\label{met:rs:universe-symbology-and-corporate-actions:buffer}
At each rebalance date, rank the candidates by the liquidity measure known that day. A member keeps its place while its rank
is within an exit rank wider than the target size (say 600 for a universe of 500); the free places go to the best-ranked
non-members. Index providers use the same device (buffer rules, One Quant Book 1, chapter 15).
\end{method}

In the chapter's simulated market, a universe of the 500 largest names by traded value replaces 11.5\% of its members each
month without a buffer and 3.1\% with an exit rank of 600 (\cref{fig:rs:universe-symbology-and-corporate-actions:churn}). The
buffered universe is not less liquid in any way that matters: its members are all within the top 600.

\begin{omfigure}
\begin{tikzpicture}
\begin{axis}[width=11cm,height=4.8cm,xlabel={month},ylabel={members replaced},tick label style={font=\small},
  label style={font=\small},xmin=0,xmax=120,ymin=0,ymax=100,grid=major,grid style={gray!25},legend columns=2,
  legend style={font=\footnotesize,at={(0.5,-0.32)},anchor=north,draw=none,fill=none,
  /tikz/every even column/.append style={column sep=8pt}},table/col sep=comma]
\addplot[omThm,thick] table[x=month,y=plain]{figdata/research/04-universe-symbology-and-corporate-actions/churn.csv};
\addplot[omDef,very thick] table[x=month,y=buffer]{figdata/research/04-universe-symbology-and-corporate-actions/churn.csv};
\legend{top 500 each month,top 500 with an exit rank of 600}
\end{axis}
\end{tikzpicture}

\omcaption{Members of a 500-name liquidity universe replaced at each monthly rebalance, in the chapter's simulated market:
11.5\% a month on average without a buffer, 3.1\% with one. Data: the chapter's module, seeded.}\label{fig:rs:universe-symbology-and-corporate-actions:churn}
\end{omfigure}

\section{Identifiers that change}

One Quant Book 1, chapter 28, made the rule: a ticker is not an identifier. Tickers are changed when a company renames
itself, freed when it delists, and given to a new company later. The example is not rare. After General Motors Corporation
filed for Chapter 11 on 1 June 2009, its shares traded over the counter (as MTLQQ by the end of that year) and the ticker GM
was free; the new General Motors Company, which had bought most of the old company's assets, listed on the New York Stock
Exchange under GM in November 2010. A table keyed on GM holds two different companies, separated by a bankruptcy.

National identifiers are more stable but not permanent. A CUSIP is a 9-character code for US and Canadian securities, and an
ISIN (ISO 6166) a 12-character code made of a two-letter country prefix, a national identifier and a check digit; both
identify an issue, not a company, and belong to the external symbols a firm must map rather than trust. The Financial Instrument Global Identifier
(FIGI) is designed not to change: once assigned it is never changed, and when the instrument ceases to exist the FIGI is
retired and never reused.

\begin{definition}[Permanent identifier, identifier mapping, ticker change]\label{def:rs:universe-symbology-and-corporate-actions:permanent}
A \emph{permanent identifier}\index{permanent identifier} is an identifier assigned to a security when it enters the firm's
data and never changed or reused, whatever happens to its external symbols. An \emph{identifier mapping}\index{identifier mapping}
is the table that links each external symbol (ticker, CUSIP, ISIN, vendor code) to a permanent identifier over a date range.
A \emph{ticker change}\index{ticker change} ends one row of that mapping and starts another for the same permanent identifier.
\end{definition}

\begin{omfigure}
\begin{tikzpicture}[font=\footnotesize,x=1cm,y=0.9cm]
\draw[->,>=Stealth] (0,0) -- (12.6,0) node[right] {time};
\foreach \x/\lab in {0/2008,2/2009,4/2010,6/2011,8/2012,10/2013,12/2014} \draw (\x,0.08) -- (\x,-0.08) node[below] {\lab};
\draw[omDef,line width=5pt] (0,3.0) -- (2.8,3.0);
\node[omDef,anchor=south west] at (0,3.25) {permanent id 1044: ticker GM};
\draw[omDef!50,line width=5pt] (2.8,3.0) -- (4.2,3.0);
\node[omDef!70,anchor=west] at (4.35,3.0) {1044: over the counter};
\draw[omThm,line width=5pt] (5.8,1.6) -- (12.4,1.6);
\node[omThm,anchor=south west] at (5.8,1.85) {permanent id 2877: ticker GM};
\draw[black!60,dashed] (0,0.75) -- (2.8,0.75);
\draw[black!60,dashed] (5.8,0.75) -- (12.4,0.75);
\node[black!70,anchor=south west] at (6.9,0.8) {a table keyed on GM: two companies, one series};
\draw[->,>=Stealth,black!60] (2.8,0.75) to[out=20,in=160] (5.8,0.75);
\end{tikzpicture}

\omcaption{A ticker reused. The listing with \omterm{def:rs:universe-symbology-and-corporate-actions:permanent}{permanent identifier} 1044 trades as GM until mid-2009, then over the counter under
another symbol; from late 2010 a new listing, 2877, trades as GM. A \omterm{def:rs:universe-symbology-and-corporate-actions:master}{security master} keeps two identifiers and two date ranges;
a table keyed on the ticker joins them into one series with a gap. The identifiers are illustrative.}\label{fig:rs:universe-symbology-and-corporate-actions:reuse}
\end{omfigure}

\section{The security master}

\begin{definition}[Security master]\label{def:rs:universe-symbology-and-corporate-actions:master}
A \emph{security master}\index{security master} is the firm's reference table of securities: for each \omterm{def:rs:universe-symbology-and-corporate-actions:permanent}{permanent identifier}, its
issuer, its listing and delisting dates, its external identifiers over time, its corporate actions and its delisting reason and
return; every other table refers to securities by \omterm{def:rs:universe-symbology-and-corporate-actions:permanent}{permanent identifier} only.
\end{definition}

The \omterm{def:rs:universe-symbology-and-corporate-actions:master}{security master} answers, as of any date, three questions every research table asks: which security did this ticker mean
on that date, what was this security's ticker then, and what was listed. It is built from events, not from a snapshot: a listing
creates an identifier, a rename ends a mapping and starts another, a delisting closes the security and records how it ended,
and a split or a dividend becomes an adjustment factor (Book 1, chapter 8) attached to the identifier, with its announcement and
ex-dates.

The simulated market shows what joining on tickers does. Over ten years it has 1\,471 listings, of which 471 end: 113 for
performance (a $-30\%$ \omterm{def:rs:universe-symbology-and-corporate-actions:delisting}{delisting return}) and 358 in cash mergers (at a premium); 158 tickers are reused. A price file keyed on
tickers therefore contains 164 returns that run from one company's last price to another's first, with an average absolute
size of 164\%. Ticker KMKN is an example (\cref{fig:rs:universe-symbology-and-corporate-actions:splice}): its first holder last
traded at 4.23 before its delisting; three months later a new company listed under KMKN at 39.55. In the ticker file, KMKN
``rose'' 835\% in three months.

\begin{omfigure}
\begin{tikzpicture}
\begin{axis}[width=11cm,height=4.8cm,xlabel={month},ylabel={price},tick label style={font=\small},
  label style={font=\small},xmin=0,xmax=120,ymin=0,grid=major,grid style={gray!25},legend columns=2,
  legend style={font=\footnotesize,at={(0.5,-0.32)},anchor=north,draw=none,fill=none,
  /tikz/every even column/.append style={column sep=8pt}},table/col sep=comma,unbounded coords=jump]
\addplot[omDef,very thick] table[x=month,y=old]{figdata/research/04-universe-symbology-and-corporate-actions/splice.csv};
\addplot[omThm,very thick] table[x=month,y=new]{figdata/research/04-universe-symbology-and-corporate-actions/splice.csv};
\legend{first holder of KMKN,second holder of KMKN}
\end{axis}
\end{tikzpicture}

\omcaption{The ticker KMKN in the simulated market's ticker-keyed price file: the first holder falls to 4.23 and is delisted
in month 36; a new company lists under the same ticker in month 39 at 39.55. A table keyed on the ticker records a return of
$+835\%$ between the two. Data: the chapter's module, seeded.}\label{fig:rs:universe-symbology-and-corporate-actions:splice}
\end{omfigure}

\section{Adjusting without leaking}

Prices are adjusted so that splits and dividends do not look like returns: the adjusted series is the traded price divided by
the product of the adjustment factors of all later actions (Book 1, chapter 8). The adjustment is exact for returns and wrong for
levels, and the difference leaks.

\begin{proposition}[What adjusted prices may and may not be used for]\label{prop:rs:universe-symbology-and-corporate-actions:adjusted}
Returns computed from an adjusted series are the total returns of holding the security. A level of the adjusted series at a
past date, however, depends on the actions that happened \emph{after} that date: a rule applied to it at date $t$ uses information
from after $t$.
\end{proposition}

\begin{proof}
Between two actions the adjustment divides all prices by the same constant, which cancels in a ratio; at an action's ex-date the
factor removes the mechanical price change. The level at $t$ is the traded price divided by the product of the factors of the
actions after $t$, which are not known at $t$.
\end{proof}

A price floor is the classic case. A research universe excludes stocks under USD~5, and the floor is applied to the adjusted
prices of a vendor file. A stock that later splits two-for-one had, in the adjusted file, half its traded price: at 8.65 it traded
comfortably above the floor, but its adjusted price, 4.32, sits below it
(\cref{fig:rs:universe-symbology-and-corporate-actions:floor}). Why did it split? Because its price rose. The floor on adjusted
prices therefore excludes, from the early universe, stocks selected for their future rise. In the simulated market it wrongly
drops 0.3\% of the name-months above the floor, and the stocks it drops return 71\% over the next twelve months, against 15\%
for the universe. A strategy that looks good at buying low-priced stocks, or bad at it, may be measuring this leak.

\begin{omfigure}
\begin{tikzpicture}
\begin{semilogyaxis}[width=11cm,height=5cm,xlabel={month},ylabel={price (log scale)},tick label style={font=\small},
  label style={font=\small},xmin=0,xmax=120,ymin=2,ymax=300,ytick={2,5,10,20,50,100,200},
  yticklabels={2,5,10,20,50,100,200},grid=major,grid style={gray!25},legend columns=3,
  legend style={font=\footnotesize,at={(0.5,-0.32)},anchor=north,draw=none,fill=none,
  /tikz/every even column/.append style={column sep=8pt}},table/col sep=comma]
\addplot[omDef,very thick] table[x=month,y=traded]{figdata/research/04-universe-symbology-and-corporate-actions/adjfloor.csv};
\addplot[omThm,very thick,dashed] table[x=month,y=adjusted]{figdata/research/04-universe-symbology-and-corporate-actions/adjfloor.csv};
\addplot[black!70,thick,sharp plot] coordinates {(0,5) (120,5)};
\legend{traded price,adjusted for the later split,USD 5 floor}
\end{semilogyaxis}
\end{tikzpicture}

\omcaption{A simulated stock that splits two-for-one in month 57. Before the split its traded price is above USD~5, but its
split-adjusted price, which divides by the later factor, starts below it: a floor on adjusted prices drops the stock for having
risen afterwards. Data: the chapter's module, seeded.}\label{fig:rs:universe-symbology-and-corporate-actions:floor}
\end{omfigure}

The rule that follows is simple. Returns from adjusted prices; every rule on a level (a price floor, a capitalisation cut-off, a
distance from a past high) on the price that traded that day, with the shares outstanding known that day.

\section{Delistings and the delisting return}

\begin{definition}[Delisting return]\label{def:rs:universe-symbology-and-corporate-actions:delisting}
The \emph{delisting return}\index{delisting return} of a security is the return from its last regular trading price to the value
its holders actually received or could realise: the cash of a merger, the price of the shares on the market where they continued
to trade, or what a bankruptcy left.
\end{definition}

The \omterm{def:rs:universe-symbology-and-corporate-actions:delisting}{delisting return} is the last return of a security's history, and it is the one most often missing. For a cash merger it is
known and usually small (the price had converged to the deal). For a performance delisting it is large and negative: Shumway
(1997) found the returns of NYSE and AMEX stocks delisted for performance reasons mostly missing from the standard US research
database, and, recovering 71\% of them from over-the-counter prices, an average of $-30\%$. A backtest that lets the position
vanish at its last regular price has treated $-30\%$ as zero, for exactly the positions that were losing. With a \omterm{def:rs:universe-symbology-and-corporate-actions:master}{security master}
the rule is mechanical: every delisting has a reason and a return; if the return is not known, an estimate by reason is used and
flagged, never zero.

\section{Tutorial: a market with changing names}\label{tut:rs:universe-symbology-and-corporate-actions:master}

\begin{tutorial}
\textbf{Goal.} Simulate ten years of listings, build their \omterm{def:rs:universe-symbology-and-corporate-actions:master}{security master} and universes, and measure what tickers and adjusted
prices do to a research table. \textbf{End state:} \cref{fig:rs:universe-symbology-and-corporate-actions:churn,fig:rs:universe-symbology-and-corporate-actions:splice,fig:rs:universe-symbology-and-corporate-actions:floor};
164 spliced returns; churn 11.5\% and 3.1\%; the dropped names' 71\% against 15\%.
\begin{tutsteps}
\item \textbf{Resolve as of a date.} A ticker belongs to one \omterm{def:rs:universe-symbology-and-corporate-actions:permanent}{permanent identifier} on a date; a listing has one ticker on a date,
or none before its listing and after its delisting.
\omcode{code/firm/secmaster/firm_secmaster.py}{73}{89}{As-of resolution in both directions.}\label{lst:rs:universe-symbology-and-corporate-actions:resolve}
\item \textbf{Universes.} Rank by the traded value known each month; keep members within the exit rank; fill with the best
non-members.
\omcode{code/firm/secmaster/firm_secmaster.py}{107}{124}{A buffered membership rule.}\label{lst:rs:universe-symbology-and-corporate-actions:buffer}
\item \textbf{The two leaks.} Build the ticker-keyed price file with \texttt{ticker\_panel} and count the returns that span two
listings; apply the USD~5 floor to adjusted and to traded prices with \texttt{price\_filter\_leak}.
\end{tutsteps}
\textbf{What to change next.} Set \texttt{reuse\_prob=0} and check that the spliced returns disappear; add a minimum history of
twelve months to the universe rule and measure how many new listings it delays.
\end{tutorial}

\section{Build: the security master}\label{bld:rs:universe-symbology-and-corporate-actions:secmaster}

\begin{build}
\bfield{Purpose} The miniature firm's one reference table: every price, feature, position and trade in later chapters refers to a
\omterm{def:rs:universe-symbology-and-corporate-actions:permanent}{permanent identifier}, and every symbol is resolved through here as of the date of use. The synthetic universe of chapter~5 and
the backtesters of Part~IV use it.
\bfield{Interface} \texttt{SecurityMaster().list(pid, ticker, date)}, \texttt{rename(pid, ticker, date)}, \texttt{delist(pid, date,
reason, ret)}; \texttt{resolve(ticker, date)}, \texttt{ticker(pid, date)}, \texttt{listed(date)}, \texttt{spans(pid)},
\texttt{delisting(pid)}, \texttt{reused()}; \texttt{buffered\_universe(dates, scores, n, exit\_rank)}; \texttt{churn(universe)}.
\bfield{Rules} A \omterm{def:rs:universe-symbology-and-corporate-actions:permanent}{permanent identifier} is used once; a ticker is held by at most one listing on any date; a rename takes effect on its
date (the old ticker is not held that day), a delisting on its date is the last day held; universes read only the scores of their
date.
\bfield{Acceptance tests} \texttt{code/firm/secmaster/tests/}: a \omterm{def:rs:universe-symbology-and-corporate-actions:permanent}{ticker change} resolved on both sides of its date; a reused ticker
resolved to two listings, with the delisting recorded; a held ticker refused; a buffer cuts churn.
\bfield{Stretch} Corporate actions attached to identifiers with announcement dates (wrapping Book 1's \texttt{firm.corpactions});
several external identifier types with their own date ranges; mergers that map one listing into another (a stock merger).
\end{build}

\begin{omsources}
\item Meta Platforms, Inc., ``Meta Platforms, Inc. to change ticker symbol to `META' on June 9'', press release, 31 May 2022 (SEC
filing, exhibit 99.1).
\item General Motors Company, Form S-1/A registration statement, November 2010 (listing under ``GM''; Motors Liquidation Company
shares trading as MTLQQ).
\item ISO TC 68, ``What is ISIN?''; CUSIP Global Services, identifier descriptions; OpenFIGI, ``About FIGI''.
\item T.~Shumway, ``The delisting bias in CRSP data'', \emph{Journal of Finance} 52(1), 1997.
\end{omsources}

\section{Exercises}

\begin{exercise}[$\star$]\label{exo:rs:universe-symbology-and-corporate-actions:1}
A stock trades at USD~12 and splits three-for-one a year later. What is its adjusted price today for that date, and does a USD~5
floor on adjusted prices include it?
\end{exercise}

\begin{exercise}[$\star$]\label{exo:rs:universe-symbology-and-corporate-actions:2}
A table keyed on tickers holds a company that was delisted at USD~0.80 and, two months later, another listed under the same ticker
at USD~24. What return does the table show?
\end{exercise}

\begin{exercise}[$\star$]\label{exo:rs:universe-symbology-and-corporate-actions:3}
An equally weighted universe of 500 names replaces 11.5\% of its members a month; each replaced name is sold and its replacement
bought, at 20 basis points of the traded value on each side. What does membership churn alone cost a year, as a share of capital?
And with the buffer's 3.1\%?
\end{exercise}

\begin{exercise}[$\star\star$]\label{exo:rs:universe-symbology-and-corporate-actions:4}
In a year, 2\% of a universe's names are delisted for performance with no \omterm{def:rs:universe-symbology-and-corporate-actions:delisting}{delisting return} in the data, and a strategy holds them
at average weight. By how much does ignoring a $-30\%$ \omterm{def:rs:universe-symbology-and-corporate-actions:delisting}{delisting return} overstate the equal-weighted return?
\end{exercise}

\begin{exercise}[$\star\star$]\label{exo:rs:universe-symbology-and-corporate-actions:5}
Why does an exit rank wider than the target size lower churn more than it lowers liquidity? Give an argument in terms of how ranks
move from month to month.
\end{exercise}

\begin{exercise}[$\star\star$]\label{exo:rs:universe-symbology-and-corporate-actions:6}
A signal uses the distance of a price from its 52-week high. On which series should it be computed: traded, split-adjusted or
total-return adjusted prices? What goes wrong with each of the others?
\end{exercise}

\begin{exercise}[$\star\star\star$]\label{exo:rs:universe-symbology-and-corporate-actions:7}
\emph{Coding.} Rerun the simulated market with \texttt{reuse\_prob} 0, 0.35 and 0.7. How do the number of spliced returns and
their mean and median absolute size change? Why is the mean so much larger than the median?
\end{exercise}

\begin{exercise}[$\star\star\star$]\label{exo:rs:universe-symbology-and-corporate-actions:8}
\emph{Find the flaw.} ``Our universe is the stocks in our vendor's current file with an adjusted price above USD~5 and a history of
at least three years. We rebuild it every month back to 2010 from the same file.''
\end{exercise}

\section{Problem: The Ticker That Changed Hands}

\begin{problem}[{Weekend problem --- what a ticker-keyed table and an adjusted price floor do to a universe}]\label{pb:rs:universe-symbology-and-corporate-actions:1}
The chapter's simulated market: 1\,000 listings at any time over 120 months, performance delistings at a cumulative log return
of $-1.6$ with $-30\%$, cash mergers at 0.3\% a month with a 25\% premium, splits above USD~200, renames at 0.2\% a month, and
freed tickers reused by 35\% of new listings, seed 4.

\medskip\noindent\textbf{Part I --- The market.}
\begin{enumerate}
\item How many listings are there over the ten years, and how many end?
\item How many end for performance and how many by merger?
\item How many splits are there, in how many names?
\item How many tickers are reused, out of how many tickers ever used?
\end{enumerate}

\medskip\noindent\textbf{Part II --- The ticker-keyed file.}
\begin{enumerate}[resume]
\item How many returns in the ticker file span two listings?
\item What is their average absolute size?
\item What return does the file show for KMKN between months 36 and 39?
\item What share of the (month, ticker) pairs of the buffered universe have a twelve-month history mixing two listings?
\item Why is that share much smaller than the share of reused tickers?
\end{enumerate}

\medskip\noindent\textbf{Part III --- The adjusted floor.}
\begin{enumerate}[resume]
\item What share of the name-months above USD~5 does a floor on adjusted prices wrongly drop?
\item What is the next-twelve-month return of the dropped names, and of the universe?
\item Why are the dropped names future winners?
\item What is the example stock's traded and adjusted price at month 0, and when does it split?
\item Which other level-based rules share the defect?
\end{enumerate}

\medskip\noindent\textbf{Part IV --- Universes.}
\begin{enumerate}[resume]
\item What are the monthly churns of the 500-name universe without and with an exit rank of 600?
\item What does the buffer cost in liquidity?
\item Where do the delisted names' last returns come from in the simulation, and in real data?
\item State the \emph{named result}: the number and average size of the spliced returns, and the forward return of the names the
adjusted floor drops against the universe's.
\item Which two tables does a \omterm{def:rs:universe-symbology-and-corporate-actions:master}{security master} make unnecessary?
\item In one sentence: why must a research table never be keyed on a ticker?
\end{enumerate}
\end{problem}

\section{Interview questions}

\begin{interviewq}[$\star$ \iqroles{developer, researcher}]\label{iq:rs:universe-symbology-and-corporate-actions:1}
Why is a ticker not a good key for a research table? What would you use instead?
\end{interviewq}

\begin{interviewq}[$\star\star$ \iqroles{researcher}]\label{iq:rs:universe-symbology-and-corporate-actions:2}
You filter your universe with a USD~5 minimum price on adjusted data. What bias did you just introduce?
\end{interviewq}

\begin{interviewq}[$\star\star$ \iqroles{developer}]\label{iq:rs:universe-symbology-and-corporate-actions:3}
Design the schema of a \omterm{def:rs:universe-symbology-and-corporate-actions:master}{security master} that answers ``which security had ticker X on date D?'' efficiently.
\end{interviewq}

\begin{interviewq}[$\star\star$ \iqroles{researcher, trader}]\label{iq:rs:universe-symbology-and-corporate-actions:4}
How do you choose the size and the rebalancing rule of a stock universe for a daily strategy?
\end{interviewq}

\begin{interviewq}[$\star\star$ \iqroles{researcher, risk}]\label{iq:rs:universe-symbology-and-corporate-actions:5}
Your backtest holds a stock that is delisted after a bankruptcy filing. What return do you book, and why?
\end{interviewq}

\begin{interviewq}[$\star\star\star$ \iqroles{researcher}]\label{iq:rs:universe-symbology-and-corporate-actions:6}
Prove that returns computed from split-adjusted prices equal holding returns, and show a rule on adjusted levels that looks into the
future.
\end{interviewq}
